\begin{theorem}[Existence of exact character conductors]
\label{mod:basic-characterconductorexistence}
Let $F$ be a nonarchimedean local field.

For every continuous quasi-character
$\chi:F^\times\to\mathbf C^\times$, there is an $m\in\mathbf N$
such that
\[
 \operatorname{IsMultiplicativeConductor}(F,\chi,m).
\]
Thus $U_F^m$ is the first unit-filtration layer on which $\chi$ is trivial,
including the endpoint convention that $m=0$ means triviality on
$U_F^0=\mathcal O_F^\times$ rather than global triviality.

For every nontrivial continuous additive character
$\psi:F\to\mathbf C^\times$, there is an $n\in\mathbf Z$ such
that
\[
 \operatorname{IsAdditiveConductor}(F,\psi,n).
\]
Equivalently, $\mathfrak p_F^{-n}$ is the largest fractional valuation
lattice on which $\psi$ is trivial.  The nontriviality hypothesis is essential
for the lower boundary and must not be weakened or suppressed.

For the multiplicative assertion, first map every nonzero complex value to
its unit-circle phase.  Continuity and the positive unit-filtration
neighborhood basis put some $U_F^{k+1}$ inside the centered arc of radius
$\pi/2$.  Since this subgroup is closed under positive powers, circle
small-subgroup rigidity makes the phase trivial there.  Compact-unit
unitarity then makes $\chi$ itself trivial there.  The least such depth,
chosen by natural-number well-ordering, is the exact conductor.

The additive proof is analogous but uses a compact fractional lattice.
Boundedness of its character image proves unitarity on that lattice, while
the lattice neighborhood basis and circle rigidity prove triviality on some
sufficiently deep lattice $\mathfrak p_F^r$.  If $\psi\ne1$, choose
$x$ with $\psi(x)\ne1$.  Writing $\operatorname{ord}_F(x)=k$, every trivial
depth is at least $k+1$; hence the trivial depths are bounded below.  The
least-integer theorem produces their least member $r$, and the additive
conductor in the manuscript's convention is $-r$.

Existence and uniqueness define canonical conductor exponents and canonical
constructors
\[
 \chi\longmapsto\texttt{LocalQuasiCharData}(F),\qquad
 (\psi,\psi\ne1)\longmapsto\texttt{LocalAddCharData}(F),
\]
whose character fields are definitionally the original characters.  These
packages are the unique exact-conductor data having their stated characters.
In particular, proof-bearing additive or multiplicative character data are
equal whenever their character fields are equal.  The additive constructor
retains an explicit proof that the input character is nontrivial.

The continuity arguments use the public neighborhood-basis theorems for the
fractional lattices and unit filtration.  No private replacement for those
foundational cofinality statements is introduced in this node.
\LeanFile{LanglandsFirstMainLemma/Basic/CharacterConductorExistence.lean}
\lean{LanglandsFirstMainLemma.continuousAddChar_norm_eq_one_of_mem_lattice}
\lean{LanglandsFirstMainLemma.exists_quasiCharTrivialOnUnitFiltration}
\lean{LanglandsFirstMainLemma.exists_addCharTrivialOnLattice}
\lean{LanglandsFirstMainLemma.exists_isMultiplicativeConductor}
\lean{LanglandsFirstMainLemma.exists_isAdditiveConductor}
\lean{LanglandsFirstMainLemma.multiplicativeConductorExponent}
\lean{LanglandsFirstMainLemma.additiveConductorExponent}
\lean{LanglandsFirstMainLemma.LocalQuasiCharData.eq_of_character_eq}
\lean{LanglandsFirstMainLemma.LocalAddCharData.eq_of_character_eq}
\lean{LanglandsFirstMainLemma.canonicalLocalQuasiCharData}
\lean{LanglandsFirstMainLemma.canonicalLocalQuasiCharData_character}
\lean{LanglandsFirstMainLemma.canonicalLocalAddCharData}
\lean{LanglandsFirstMainLemma.canonicalLocalAddCharData_character}
\lean{LanglandsFirstMainLemma.characterConductorExistence}
\leanok
\SourceText{Definitions Multiplicative conductor and Additive conductor; local notation, subsection Local fields, ideals, and characters.}
\uses{mod:basic-characterconductors,mod:localfield-lattices,mod:localfield-unitfiltration}
\FormalizationContract
The Lean theorem \texttt{characterConductorExistence} is the conjunction of
the two quantified existence assertions above.  The accompanying declarations
prove eventual triviality from continuity, construct exact exponents, and
package the canonical data.  No trivial additive character is packaged.
\end{theorem}
